%%%%%%%%%%%%%%%%%%%%%%%%%%%%%%%%%%%%%%%%%%%%%%%%%%
%                                                %
%  Use latex to format this document.            %
%  AGASHE - RIBET - STEIN : Congruences          %
%                                                %
%%%%%%%%%%%%%%%%%%%%%%%%%%%%%%%%%%%%%%%%%%%%%%%%%%

%
%    Documenta-LaTeX.sty (2e)
%
%    Version 0.5
%
%    Copyright (C) 2000, 2001 Ulf Rehmann
%    rehmann@mathematik.uni-bielefeld.de
%
% ------------------------------ PROLOGUE -----------------------------
%
% This prologue is to be prepended to your LaTeX file in order to
% give it the layout required by Documenta Mathematica.
% Please edit the \documentclass line below, add \usepackage lines
% according to your needs and
% go to the line starting with \Title below, fill in title,
% authors' names and so on and follow the further instructions there.
% 
% Choose one of the commands \documentclass[]{} or \documentstyle[]{}
% depending on the version of LaTeX you use:
%
%
% If you use LaTeX 2e use the following line (or similar):
\documentclass[10pt,twoside]{article}

%\usepackage[hyperindex,pdfmark]{hyperref}
\usepackage[hypertex]{hyperref}
%\usepackage[active]{srcltx}

\newcommand{\edit}[1]{\footnote{EDIT: #1}\marginpar{\hfill {\sf \thefootnote}}}
\usepackage{amsmath}
\usepackage{amsfonts}
\usepackage{amssymb}
\usepackage{amsthm}


%*************Declaration Section******************
\newcommand{\Z}{{\bf{Z}}}
\newcommand{\Q}{{\bf{Q}}}
\newcommand{\Qbar}{{\overline{\Q}}}
\newcommand{\Zbar}{{\overline{\Z}}}
\newcommand{\R}{{\bf{R}}}
\newcommand{\C}{{\bf{C}}}
\newcommand{\F}{{\bf{F}}}
\newcommand{\T}{{\bf{T}}}
\newcommand{\Tf}{\T_{\scriptscriptstyle [f]}}
\newcommand{\J}{\bf{J}}
\newcommand{\m}{{\mathfrak{m}}}
\newcommand{\PP}{\bf{P}}
\newcommand{\RP}{{\noindent \bf Research problems:}}
\newcommand{\B}{{\noindent $\bullet$ \ }}
%\newcommand{\HH}{\rm{H}}
%\newcommand{\SS}{\rm{S}}
%\newcommand{\LL}{\rm{L}}
\newcommand{\dimq}{{\dim_{\scriptscriptstyle \Q}}}
\newcommand{\tensor}{\otimes}
\newcommand{\ten}{\otimes}
\newcommand{\ra}{\rightarrow}
\newcommand{\lra}{\longrightarrow}
\newcommand{\ua}{\uparrow}
\newcommand{\da}{\downarrow}
\newcommand{\hra}{\hookrightarrow}
\newcommand{\lan}{\langle}
\newcommand{\ran}{\rangle}
\font\cyr=wncyr10 scaled \magstep 1
%\font\cyr=wncyr10p
\newcommand{\Sha}{\hbox{\cyr X}}
\newcommand{\phie}{\phi_{\scriptscriptstyle E}}
\newcommand{\phia}{\phi_{\scriptscriptstyle A}}
\newcommand{\vphi}{\varphi}
\newcommand{\dphif}{{\rm{deg}} \phi_{\scriptscriptstyle F}}
\newcommand{\re}{r_{\scriptscriptstyle E}}
\newcommand{\me}{m_{\scriptscriptstyle E}}
\newcommand{\qe}{q{\mbox{\rm-exp}}}
\newcommand{\Fe}{F{\mbox{\rm-exp}}}
\newcommand{\sqe}{{\scriptstyle{q}}{\mbox{\scriptsize{\rm-exp}}}}
\newcommand{\sFe}{{\scriptstyle{F}}{\mbox{\scriptsize{\rm-exp}}}}
%\newcommand{\Zm}{{\Z}\left[\frac{1}{m}\right]}
\newcommand{\Zm}{{\Z}[1/m]}
\newcommand{\tZm}{{\textstyle \Zm}}
\newcommand{\HM}{H^0(M_0(N)_S,\Omega)}
\newcommand{\HJ}{H^0(J_S, \Omega_{J/S}^1)}
\newcommand{\GS}[2]{H^0(#1_#2, \Omega_{#1/#2}^1)}
\newcommand{\po}{\phi_1}
\newcommand{\pt}{\phi_2}
\newcommand{\Fp}{{{\F}_p}}
\newcommand{\tfp}{\tensor \Fp}
\newcommand{\HH}{\rm{H}}
\newcommand{\SSS}{\rm{S}}
\newcommand{\nerA}{{\mathcal{A}}}
\newcommand{\ce}{c_{\scriptscriptstyle{E}}}
\newcommand{\cje}{c_{\scriptscriptstyle{J_e}}}
\newcommand{\ca}{{c_{\scriptscriptstyle{A}}}}
\newcommand{\caf}{{c_{\scriptscriptstyle{A_f}}}}
\newcommand{\na}{{n_{\scriptscriptstyle{A}}}}
\newcommand{\ma}{{m_{\scriptscriptstyle{A}}}}
\newcommand{\nje}{{n_{\scriptscriptstyle{J_e}}}}
\newcommand{\nel}{{n_{\scriptscriptstyle{E}}}}
\newcommand{\rA}{{r_{\scriptscriptstyle{A}}}}
\newcommand{\rAe}{{\tilde{r}_{\scriptscriptstyle{A}}}}
\newcommand{\rAfe}{{\tilde{r}_{\scriptscriptstyle{A_f}}}}
\newcommand{\rAf}{{r_{\scriptscriptstyle{A_f}}}}
\newcommand{\mA}{{m_{\scriptscriptstyle{A}}}}
\newcommand{\nA}{{n_{\scriptscriptstyle{A}}}}
\newcommand{\nAe}{{\tilde{n}_{\scriptscriptstyle{A}}}}
\newcommand{\nAfe}{{\tilde{n}_{\scriptscriptstyle{A_f}}}}
\newcommand{\nAf}{n_{\scriptscriptstyle{A_f}}}
\newcommand{\If}{{I_{\scriptscriptstyle{f}}}}
\newcommand{\tauc}{\tau_{\scriptscriptstyle{\C}}}
\newcommand{\piAd}{\pi_{\scriptscriptstyle{\Adual}}}
\newcommand{\piB}{\pi_{\scriptscriptstyle{B}}}
\newcommand{\Adual}{A^{\vee}}
\newcommand{\Bdual}{B^{\vee}}
\newcommand{\Edual}{E^{\vee}}
\newcommand{\Iehat}{I_e^{\vee}}
\newcommand{\Jdual}{J^{\vee}}
\newcommand{\Jehat}{J_e^{\vee}}
\newcommand{\Hehat}{H_e^{\vee}}
\newcommand{\Hje}{{H_1(J_e,\Z)}}
\newcommand{\U}{{\bf U}}
\newcommand{\cusps}{{\rm cusps}}
\newcommand{\Hom}{{\rm Hom}}
\newcommand{\OO}{{\mathcal{O}}}
\newcommand{\EJ}{{{\rm End}_{\scriptscriptstyle \Q} J_0(N)}}
\newcommand{\Tzp}{{\T \tensor \Z_p}}
\newcommand{\Tm}{{\T_{(\m)}}}
\newcommand{\Mm}{{M_{(\m)}}}
\newcommand{\Tpm}{{\T'_{(\m)}}}
\newcommand{\End}{{\rm End}}
\newcommand{\Endq}{{\End^0}}
\newcommand{\cA}{\mathcal{A}}
\newcommand{\cB}{\mathcal{B}}
\newcommand{\cJ}{\mathcal{J}}
\newcommand{\cX}{\mathcal{X}}
\newcommand{\xra}{\xrightarrow}
\newcommand{\isom}{\cong}
\newcommand{\Fell}{\mathbf{F}_\ell}
\newcommand{\Zell}{\mathbf{Z}_\ell}
\newcommand{\Tor}{\mbox{\rm Tor}}
\newcommand{\invlim}{\rm invlim}
\newcommand{\Xmu}{\mathcal{X}_{\mu}}
\renewcommand{\H}{\HHH}

\newcommand{\comment}[1]{}



\newtheorem{lem}{Lemma}[section]
\newtheorem{cor}[lem]{Corollary}
\newtheorem{prop}[lem]{Proposition}
\newtheorem{conj}[lem]{Conjecture}
\newtheorem{thm}[lem]{Theorem}

\theoremstyle{definition}
\newtheorem{defi}[lem]{Definition}
\newtheorem{rmk}[lem]{Remark}
\newtheorem{que}[lem]{Question}
\newtheorem{eg}[lem]{Example}


\DeclareMathOperator{\Ann}{Ann}
\DeclareMathOperator{\sss}{ss}
\renewcommand{\ss}{\sss}
\DeclareMathOperator{\Ext}{Ext}
\DeclareMathOperator{\reg}{reg}
\DeclareMathOperator{\ord}{ord}
\DeclareMathOperator{\Jac}{Jac}
\DeclareMathOperator{\Mat}{Mat}
\DeclareMathOperator{\Spec}{Spec}
\DeclareMathOperator{\Supp}{Supp}
\DeclareMathOperator{\Spf}{Spf}
\DeclareMathOperator{\Tan}{Tan}
\DeclareMathOperator{\Tate}{Tate}
\DeclareMathOperator{\HHH}{H}
\DeclareMathOperator{\new}{new}
\DeclareMathOperator{\old}{old}


\usepackage[all]{xy}



% If you use LaTeX 2.09 or earlier use the following line (or similar):
%\documentstyle[10pt,twoside]{article}
%
% Please go to the line starting with \Title below, fill in title,
% authors' names and so on and follow the further instructions there.
%
% The final values of the following 4 uncommented lines will be filled in
% by Doc. Math. after the manuscript has been accepted for publication.
%
\def\YEAR{\year}\newcount\VOL\VOL=\YEAR\advance\VOL by-1995
\def\firstpage{1}\def\lastpage{1000}
\def\received{}\def\revised{}
\def\communicated{}
%

\makeatletter
\def\magnification{\afterassignment\m@g\count@}
\def\m@g{\mag=\count@\hsize6.5truein\vsize8.9truein\dimen\footins8truein}
\makeatother

%%% Choose 10pt/12pt:
%\oddsidemargin1.66cm\evensidemargin1.66cm\voffset1.2cm%10pt
\oddsidemargin1.91cm\evensidemargin1.91cm\voffset1.4cm%10pt
%\magnification1200\oddsidemargin.41cm\evensidemargin.41cm\voffset-.75cm%12pt

%\textwidth12.5cm\textheight19.5cm
\textwidth12.0cm\textheight19.0cm

\font\eightrm=cmr8
\font\caps=cmcsc10                    % Theorem, Lemma etc
\font\ccaps=cmcsc10                   % Sections
\font\Caps=cmcsc10 scaled \magstep1   % Title
\font\scaps=cmcsc8

%
%-----------------headlines-----------------------------------

%\input german.sty      % uncomment if necessary
%\usepackage{german}    % uncomment if necessary

\pagestyle{myheadings}
\pagenumbering{arabic}
\setcounter{page}{\firstpage}
\def\bfseries{\normalsize\caps}

\makeatletter
\setlength\topmargin {14\p@}
\setlength\headsep   {15\p@}  
\setlength\footskip  {25\p@}  
\setlength\parindent {20\p@} 
\@specialpagefalse\headheight=8.5pt
\def\DocMath{}
\renewcommand{\@evenhead}{%
    \ifnum\thepage>\lastpage\rlap{\thepage}\hfill%
    \else\rlap{\thepage}\slshape\leftmark\hfill{\caps\SAuthor}\hfill\fi}%
\renewcommand{\@oddhead}{%
    \ifnum\thepage=\firstpage{\DocMath\hfill\llap{\thepage}}%
    \else{\slshape\rightmark}\hfill{\caps\STitle}\hfill\llap{\thepage}\fi}%
\makeatother

\def\TSkip{\bigskip}
\newbox\TheTitle{\obeylines\gdef\GetTitle #1
\ShortTitle  #2
\SubTitle    #3
\Author      #4
\ShortAuthor #5
\EndTitle
{\setbox\TheTitle=\vbox{\baselineskip=20pt\let\par=\cr\obeylines%
\halign{\centerline{\Caps##}\cr\noalign{\medskip}\cr#1\cr}}%
	\copy\TheTitle\TSkip\TSkip%
\def\next{#2}\ifx\next\empty\gdef\STitle{#1}\else\gdef\STitle{#2}\fi%
\def\next{#3}\ifx\next\empty%
    \else\setbox\TheTitle=\vbox{\baselineskip=20pt\let\par=\cr\obeylines%
    \halign{\centerline{\caps##} #3\cr}}\copy\TheTitle\TSkip\TSkip\fi%
%\setbox\TheTitle=\vbox{\let\par=\cr\obeylines%
%\halign{\centerline{\caps##} #4\cr}}\copy\TheTitle\TSkip\TSkip%
\centerline{\caps #4}\TSkip\TSkip%
\def\next{#5}\ifx\next\empty\gdef\SAuthor{#4}\else\gdef\SAuthor{#5}\fi%
\ifx\received\empty\relax
    \else\centerline{\eightrm Received: \received}\fi%
\ifx\revised\empty\TSkip%
    \else\centerline{\eightrm Revised: \revised}\TSkip\fi%
\ifx\communicated\empty\relax
    \else\centerline{\eightrm Communicated by \communicated}\fi\TSkip\TSkip%
\catcode'015=5}}\def\Title{\obeylines\GetTitle}
\def\Abstract{\begingroup\narrower
    \parskip=\medskipamount\parindent=0pt{\caps Abstract. }}
\def\EndAbstract{\par\endgroup\TSkip}

\long\def\MSC#1\EndMSC{\def\arg{#1}\ifx\arg\empty\relax\else
     {\par\narrower\noindent%
     2000 Mathematics Subject Classification: #1\par}\fi}

\long\def\KEY#1\EndKEY{\def\arg{#1}\ifx\arg\empty\relax\else
	{\par\narrower\noindent Keywords and Phrases: #1\par}\fi\TSkip}

\newbox\TheAdd\def\Addresses{\vfill\copy\TheAdd\vfill
    \ifodd\number\lastpage\vfill\eject\phantom{.}\vfill\eject\fi}
{\obeylines\gdef\GetAddress #1
\Address #2 
\Address #3
\Address #4
\EndAddress
{\def\xs{4.3truecm}\parindent=15pt
\setbox0=\vtop{{\obeylines\hsize=\xs#1\par}}\def\next{#2}
\ifx\next\empty % 1 address
     \setbox\TheAdd=\hbox to\hsize{\hfill\copy0\hfill}
\else\setbox1=\vtop{{\obeylines\hsize=\xs#2\par}}\def\next{#3}
\ifx\next\empty % 2 addresses
     \setbox\TheAdd=\hbox to\hsize{\hfill\copy0\hfill\copy1\hfill}
\else\setbox2=\vtop{{\obeylines\hsize=\xs#3\par}}\def\next{#4}
\ifx\next\empty\ % 3 addresses
     \setbox\TheAdd=\vtop{\hbox to\hsize{\hfill\copy0\hfill\copy1\hfill}
                \vskip20pt\hbox to\hsize{\hfill\copy2\hfill}}
\else\setbox3=\vtop{{\obeylines\hsize=\xs#4\par}}
     \setbox\TheAdd=\vtop{\hbox to\hsize{\hfill\copy0\hfill\copy1\hfill}
	        \vskip20pt\hbox to\hsize{\hfill\copy2\hfill\copy3\hfill}}
\fi\fi\fi\catcode'015=5}}\gdef\Address{\obeylines\GetAddress}

\hfuzz=0.1pt\tolerance=2000\emergencystretch=20pt\overfullrule=5pt
\begin{document}
%%%%% ------------- fill in your data below this line  -------------------
%%%%%    The following lines \Title ... \EndAddress must ALL be present
%%%%%    and in the given order.
\Title
%%%%%    Put here the title. Line breaks will be recognized. 
The Modular Degree, Congruence Primes and Multiplicity One
\ShortTitle 
The Modular Degree and Congruences
%%%%%    Running title for odd numbered pages, ONE line, please. 
%%%%%    If none is given, \Title will be used instead.          
\SubTitle   
%%%%%    A possible subtitle goes here.
\Author 
%%%%%    Put here name(s) of authors. Line breaks will be recognized.  
Amod Agashe\\
Kenneth A. Ribet\\
William A. Stein
\ShortAuthor 
Agashe, Ribet, Stein
%%%%%%   Running title for even numbered pages, ONE line, please. 
%%%%%%   If none is given, \Author will be used instead.          
\EndTitle
\Abstract 
%%%%%    Put here the abstract of your manuscript.
%%%%%    Avoid macros and complicated TeX expressions, as this is
%%%%%    automatically translated and posted as an html file.

The modular degree and congruence number are two fundamental
invariants of an elliptic curve over the rational field.  Frey and
M{\"u}ller have asked whether these invariants coincide.  Although
this question has a negative answer, we prove a theorem about the
relation between the two invariants: one divides the other, and the
ratio is divisible only by primes whose squares divide the conductor
of the elliptic curve.  We discuss the ratio even in the case where
the square of a prime does divide the conductor, and we study
analogues of the two invariants for modular abelian varieties of
arbitrary dimension.
 


\EndAbstract
\MSC 
%%%%%    2000 Mathematics Subject Classification: 
\EndMSC
\KEY 
%%%%%    Keywords and Phrases:     
\EndKEY
%%%%%    All 4 \Address lines below must be present. To center the last
%%%%%    entry, no empty lines must be between the following \Address
%%%%%    and \EndAddress lines.
\Address 
%%%%%    Address of first Author here
\Address 
Amod Agashe
Insert Current Address
\Address
Kenneth A. Ribet
Insert Current Address
\Address
William A. Stein
Department of Mathematics
Harvard University
Cambridge, MA  02138
{\tt was@math.harvard.edu}
\EndAddress
%%
%%       Make sure the last tex command in your manuscript
%%       before the first \end{document} is the command  \Addresses
%%
%%---------------------Here the prologue ends---------------------------------




%%--------------------Here the manuscript starts------------------------------

  
\section{Introduction}
Let~$E$ be an elliptic curve over~$\Q$.  By
\cite{breuil-conrad-diamond-taylor}, we may view~$E$ as an abelian
variety quotient over $\Q$ of the modular Jacobian $J_0(N)$, where $N$
is the conductor of~$E$.  After possibly replacing $E$ by an isogenous
curve, we may assume that the kernel of the map $J_0(N)\to E$ is
connected, i.e., that~$E$ is an {\em optimal quotient} of $J_0(N)$.

Let $f_E = \sum a_n q^n \in S_2(\Gamma_0(N))$ be the newform attached
to $E$.  The {\em congruence number}~$\re$ of~$E$ is the largest
integer such that there is an element $g =\sum b_n q^n \in
S_2(\Gamma_0(N))$ with integer Fourier coefficients $b_n$ that is
orthogonal to~$f_E$ with respect to the Peterson innner product, and
congruent to~$f_E$ modulo~$\re$ (i.e., $a_n \equiv b_n\pmod{\re}$ for
all~$n$).    The {\em modular
  degree}~$\me$ is the degree of the composite map $X_0(N)\to
J_0(N)\to E$, where we map $X_0(N)$ to $J_0(N)$ by sending $P\in
X_0(N)$ to $[P]-[\infty] \in J_0(N)$.

Section~\ref{congintro} is about relations between~$\re$ and~$\me$.
For example, $\me \mid \re$.  In \cite[Q.~4.4]{frey-muller}, Frey and
M{\"u}ller~ asked whether $\re = \me$.  We give examples in which $\re
\neq \me$, then conjecture that for any prime $p$, $\ord_p(\re/\me)
\leq \frac{1}{2}\ord_p(N)$.  We prove this conjecture when
$\ord_p(N)\leq 1$.

In Section~\ref{sec:quotients}, we consider analogues of congruence
primes and the modular degree for optimal quotients that are not
necessarily elliptic curves; these are quotients of~$J_0(N)$ and
$J_1(N)$ of any dimension associated to ideals of the relevant Hecke
algebras.  In Section~\ref{sec:main} we prove the main theorem of this
paper, and in Section~\ref{sec:mult1} we give some new examples of
failure of multiplicity one motivated by the arguments in
Section~\ref{sec:main}.

% For an introduction and the motivation for studying
% the objects in the title of the paper, the reader 
% may read Sections~\ref{sec:elliptic}
% and~\ref{sec:quotients}, skipping the proofs. 

\smallskip

{\bf \noindent Acknowledgment.} The authors are grateful to A.~Abbes,
R.~Coleman, B.~Conrad, J.~Cremona, H.~Lenstra, E.~de Shalit,
B.~Edixhoven, L.~Merel, and R.~Taylor for several discussions and
advice regarding this paper.

\section{Congruence Primes and the Modular Degree} 
\label{sec:elliptic}


Let~$N$ be a positive integer and let $X_0(N)$ be the modular curve
over~$\Q$ that classifies isomorphism classes of elliptic curves with
a cyclic subgroup of order~$N$.  The Hecke algebra~$\T$ of level~$N$
is the subring of the ring of endomorphisms of $J_0(N)=\Jac(X_0(N))$
generated by the Hecke operators $T_n$ for all $n\geq 1$.  Let~$f$ be
a newform of weight~$2$ for~$\Gamma_0(N)$ with integer Fourier
coefficients, and let $I_f$ be kernel of the 
homomorphism $\T\to \Z[\ldots, a_n(f), \ldots]$ that sends $T_n$ to
$a_n$.  Then the quotient $E = J_0(N)/I_f J_0(N)$ is an elliptic curve
over~$\Q$.  We call~$E$ the {\em optimal quotient} associated to~$f$.
Composing the embedding $X_0(N)\hra J_0(N)$ that sends $\infty$ to~$0$
with the quotient map $J_0(N) \ra E$, we obtain a surjective morphism
of curves $\phie: X_0(N) \ra E$.
\begin{defi}
The {\em modular degree} $\me$ of~$E$ is the degree of~$\phie$. 
\end{defi}

\label{congintro}

%The congruence number $\re$ and the modular degree $\me$ 
%are of great interest.  
Congruence primes have been studied by Doi, Hida, Ribet, 
Mazur and others (see, e.g.,~\cite[\S1]{ribet:modp}), 
and played an important role in Wiles's work~\cite{wiles} 
on Fermat's last theorem.  Frey and Mai-Murty have 
observed that an
appropriate asymptotic bound on the modular degree is equivalent to
the $abc$-conjecture (see~\cite[p.544]{frey:ternary} 
and~\cite[p.180]{murty:congruence}). 
Thus, results that relate congruence primes and the modular degree
are of great interest.

\begin{thm}\label{thm:ribet_au}
\label{ddivsr}
Let $E$ be an elliptic curve over $\Q$ of conductor~$N$, with modular
degree $\me$ and congruence number $\re$.  
Then $\me \mid \re$ and if $\ord_p(N)\leq 1$ then $\ord_p(\re) = \ord_p(\me)$. 
\end{thm}
We will prove a generalization of Theorem~\ref{thm:ribet_au}
in Section~\ref{sec:main} below.

The divisibility $\me\mid \re$ was first discussed
in~\cite[Th.~3]{zagier}, where it is attributed to the second author
(Ribet); however in \cite{zagier} the divisibility was mistakenly
written in the opposite direction. For some other expositions of the
proof, see~\cite[Lem~3.2]{abbull} and~\cite{cojo-kani}.  We generalize
this divisibility in Proposition~\ref{ndivsm}.  The second part of
Theorem~\ref{thm:ribet_au}, i.e., that if $\ord_p(N) \leq 1$ then
$\ord_p(\re) = \ord_p(\me)$, follows from the more general
Theorem~\ref{thm:ribet_gen} below.
%\edit{I made this change. --Amod}
%in more generality in in Section~\ref{sec:proof_ribet} below. 
Note that \cite[Prop.~3.3--3.4]{abbull} implies the weaker
statement that if $p\nmid N$ then $\ord_p(\re)=\ord_p(\me)$,
since \cite[Prop.~3.3]{abbull} 
implies $$\ord_p(\re) - \ord_p(\me) = \ord_p(\#\mathcal{C})
- \ord_p(\ce) - \ord_p(\#\mathcal{D}),$$ and by \cite[Prop.~3.4]{abbull}
$\ord_p(\#\mathcal{C}) =0$.  (Here $\ce$ is the Manin
constant of $E$, which is an integer by results of
Edixhoven and Katz-Mazur; see e.g., \cite{ars} for more details.)

Frey and M{\"u}ller~\cite[Ques.~4.4]{frey-muller} asked whether $\re =
\me$ in general.  After implementing an algorithm to compute $\re$ in
Magma \cite{magma}, we quickly found that the answer is no. The
counterexamples at conductor $N\leq 144$ are given in Table~\ref{table:moddeg}, 
where the curve
is given using the notation of \cite{cremona:alg}:
\begin{table}\caption{Differing Modular Degree and 
Congruence Number\label{table:moddeg}}
\begin{center}
\begin{tabular}{|l|l|l|}\hline
Curve & $\me$ & $\re$\\\hline
54B1 & 2 & 6\\\hline
64A1 & 2 & 4 \\\hline
72A1 & 4 & 8 \\\hline
80A1 & 4 & 8 \\\hline
88A1 & 8 & 16 \\\hline
92B1 & 6 & 12\\\hline
96A1 & 4 & 8 \\\hline
96B1 & 4 & 8 \\\hline
\end{tabular}
\begin{tabular}{|l|l|l|}\hline
Curve & $\me$ & $\re$\\\hline
99A1 & 4 & 12\\\hline
108A1 & 6 & 18\\\hline
112A1 & 8 & 16\\\hline
112B1 & 4 & 8\\\hline
112C1 & 8 & 16\\\hline
120A1 & 8 & 16\\\hline
124A1 & 6 & 12\\\hline
126A1 & 8 & 24\\\hline
\end{tabular}
\begin{tabular}{|l|l|l|}\hline
Curve & $\me$ & $\re$\\\hline
128A1 & 4 & 32\\\hline
128B1 & 8 & 32\\\hline
128C1 & 4 & 32\\\hline
128D1 & 8 & 32\\\hline
135A1 & 12 & 36\\\hline
144A1 & 4 & 8 \\\hline
144B1 & 8 & 16 \\\hline
&&\\\hline
\end{tabular}
\end{center}
\end{table}
%$$54, 64, 72, 80, 88, 92, 96, 99, 108, 112, 120, 124, 126, 128, 135,
%\text{ and } 144.$$ 

For example, the elliptic curve 54B1, given by the equation $y^2 + xy +
y = x^3 - x^2 + x - 1$, has $\re=6$ and $\me=2$.  To see explicitly
that $3 \mid \re$, observe that the newform corresponding to~$E$ is
$f=q + q^2 + q^4 - 3q^5 - q^7 + \cdots$ and the newform corresponding
to $X_0(27)$ if $g=q - 2q^4 - q^7 + \cdots$, so $g(q) + g(q^2)$
appears to be congruent to~$f$ modulo~$3$.  To prove this congruence,
we checked it for $18$ Fourier coefficients, where the 
sufficiency of precision to degree $18$
was determined using \cite{sturm:cong}.

% In accord with Theorem~\ref{thm:ribet_au}, 
%since $\ord_3(\re) \neq \ord_3(\ce)$, we have $\ord_3(54)\geq 2$. 

In our computations, there appears to be no absolute bound on the~$p$
that occur.  For example, for the curve 242B1 of conductor $N=2\cdot 11^2$
we have\footnote{The curve 242a1 in ``modern notation.''}
$$
\me = 2^4 \neq \re = 2^4\cdot 11.
$$
We propose the following replacement for Question~4.4 of
\cite{frey-muller}:
\begin{conj}\label{conj:rm}
  Let~$E$ be an optimal elliptic curve of conductor~$N$ 
  and~$p$ be any prime. 
  Then
$$
\ord_p\left(\frac{\re}{\me}\right) \leq \frac{1}{2}\ord_p(N). 
$$
\end{conj}
We verified Conjecture~\ref{conj:rm} using Magma for every optimal
elliptic curve quotient of $J_0(N)$, with $N\leq 539$.

If $p\geq 5$ then $\ord_p(N)\leq 2$, so a special case
of the conjecture is 
$$
  \ord_p\left(\frac{\re}{\me}\right) \leq 1\qquad\text{ for any }p\geq 5.
$$


\begin{rmk}
  It is often productive to parametrize elliptic curves by $X_1(N)$
  instead of $X_0(N)$ (see, e.g., \cite{stevens:param} and
  \cite{MR2135139}).  Suppose $E$ is an optimal quotient of $X_1(N)$,
  let $m_E'$ be the degree of the modular parametrization, and let
  $r_E'$ be the $\Gamma_1(N)$-congruence number, which is defined as
  above but with $S_2(\Gamma_0(N))$ replaced by $S_2(\Gamma_1(N))$.
  For the optimal quotient of $X_1(N)$ isogenous to 54B1, we find
  using Magma that $m_E' = 18$ and $r_E'=6$. Thus the equality
  $m_E'=r_E'$ fails, and the analogous divisibility $m_E'\mid
  r_E'$ no longer holds.  Also, for a curve of conductor $38$ we have
  $m_E'=18$ and $r_E'=6$, so equality need not hold even if the level
  is square free.  We hope to investigate this in a future paper.
%> N := 38; D := ND(NS(CS(ModularSymbols(Gamma1(N)))));
%> ModularDegree(D[1]); CongruenceModulus(D[1]);
%18
%6
%
%>  N := 54; D := ND(NS(CS(ModularSymbols(Gamma1(N)))));
%>  ModularDegree(D[1]); CongruenceModulus(D[1]);
%18
%18
%>  ModularDegree(D[2]); CongruenceModulus(D[2]);
%18
%6
\end{rmk}


%\subsection{Proof of Theorem~\ref{thm:ribet_au}}\label{sec:proof_ribet}

\section{Modular abelian varieties of arbitrary dimension}
\label{sec:quotients}
For $N\geq 4$, let~$\Gamma$ be a fixed choice of either~$\Gamma_0(N)$
or~$\Gamma_1(N)$, let~$X$ be the modular curve over~$\Q$ associated
to~$\Gamma$, and let~$J$ be the Jacobian of~$X$.  Let~$I$ be a {\em
  saturated} ideal of the corresponding Hecke algebra
$\T\subset\End(J)$, so $\T/ I$
is torsion free.  Then $A = A_I = J/IJ$ is an optimal quotient of~$J$
since $IJ$ is an abelian subvariety.

\begin{defi}
  If~$f=\sum a_n(f)q^n \in S_2(\Gamma)$ and $I_f=\ker(\T\to
  \Z[\ldots,a_n(f),\ldots])$, then $A=A_f=J/I_f J$ is the {\em newform
    quotient} associated to~$f$.  It is an abelian variety over~$\Q$
  of dimension equal to the degree of the field
  $\Q(\ldots,a_n(f),\ldots)$.
\end{defi}

In this section, we generalize the notions of the congruence number
and the modular degree to quotients~$A=A_I$, and state a theorem
relating the two numbers, which we prove in
Sections~\ref{sec:firstpart}--\ref{sec:secondpart}.

Let $\phi_2$ denote the quotient map $J \ra A$.  By Poincare
reducibility over $\Q$ there is a unique abelian subvariety $A^{\vee}$
of $J$ that projects isogenously to the quotient $A$ (equivalently,
which has finite intersection with $\ker(\phi_2)$), and so by Hecke
equivariance of $J \to A$ it follows that $A^{\vee}$ is $\T$-stable.
Let $\phi$ be the composite isogeny
$$
  \phi: \Adual \stackrel{\po}{\lra} J \stackrel{\pt}{\lra} A.
$$

\begin{rmk}
Note that $A^{\vee}$ is the dual abelian variety of $A$.  More
generally, if~$C$ is any abelian variety, let $C^{\vee}$ denote the
dual of~$C$.  There is a canonical principal polarization $J \cong
\Jdual$, and dualizing $\phi_2$, we obtain a map $\phi_2^\vee: \Adual
\ra \Jdual$, which we compose with $\theta^{-1}: \Jdual \cong J$ to
obtain a map $\po: \Adual \ra J$.  Note also that $\vphi$
is a polarization (induced by pullback of the theta divisor).
\end{rmk}

%\begin{prop} \label{modular:isogeny0}
%The map $\phi$ is  a polarization. 
%\end{prop}
% \begin{proof}
% Let $i$ be the injection $\phi_2^{\vee}:\Adual \ra \Jdual$, and let
% $\Theta$ denote the theta divisor.  From the definition of the
% polarization attached to an ample divisor, we see that the map~$\phi$
% is induced by the pullback $i^*(\Theta)$ of the theta divisor.  The
% theta divisor is effective, and hence so is $i^*(\Theta)$. 
%By~\cite[\S6, Application~1, p. 60]{mumford:av}, $\ker \phi$ is
%finite. Since the dimensions of $A$ and~$\Adual$ are the same, $\phi$
%is an isogeny.
% Since $\Theta$ is ample, some power of it is
% very ample. Then the pullback of this very ample power by~$i$ is again
% very ample, and hence a power of $i^*(\Theta)$ is very ample, so
% $i^*(\Theta)$ is ample (by~\cite[II.7.6]{hartshorne:ag}). 
% \end{proof}

The {\em exponent} of a finite group~$G$ is the smallest positive
integer~$n$ such that every element of~$G$ has order dividing~$n$.

\begin{defi}\label{defi:modular}
The {\em modular exponent} of~$A$ is the exponent of the kernel
of the isogeny~$\phi$, and the {\em modular number} of~$A$ is
the degree of~$\phi$. 
\end{defi}

We denote the modular exponent of~$A$ by~$\nAe$ and
the modular number by~$\nA$. 
When~$A$ is an elliptic curve, the modular
exponent is equal to the modular degree of~$A$,
and the modular number is the square of the modular degree
(see, e.g.,~\cite[p.~278]{abbull}). 
%(see \cite[p.~276]{abbull}). 
%When~$A$ is an elliptic curve, $\na$ is just the
%modular degree of~$A$. 

If~$R$ is a subring of~$\C$, 
let $S_2(R)=S_2(\Gamma;R)$ denote the subgroup of~$S_2(\Gamma)$
consisting of cups forms whose Fourier expansions at the cusp~$\infty$
have coefficients in~$R$.  (Note that $\Gamma$ is fixed for this whole
section.)
Let $S_2(\Gamma;\Z)[I]^{\perp}$ denote the orthogonal complement of
$S_2(\Gamma;\Z)[I]$ in $S_2(\Gamma;\Z)$ with respect to the Petersson inner
product. 

The following is well known, but we had difficulty finding
a good reference.
\begin{prop}
The group $S_2(\Gamma;\Z)$ is of finite rank as a $\Z$-module.
\end{prop}
\begin{proof}
  Using the standard pairing between $\T$ and $S_2(\Gamma,\Z)$ (see
  also~\cite[Theorem~2.2]{ribet:modp}) we see that $S_2(\Gamma,\Z)
  \isom \Hom(\T,\Z)$. Thus $S_2(\Gamma,\Z)$ is finitely generated
  over~$\Z$ if and only if~$\T$ is finitely generated over~$\Z$.  But
  the action of~$\T$ on $\H_1(J,\Z)$ is a faithful representation that
  embeds~$\T$ into $\Mat_{2d}(\Z) \isom \Z^{(2d)^2}$.  But~$\Z$ is
  Noetherian, so~$\T$ is finitely generated over~$\Z$.
\end{proof}

\begin{defi}\label{def:congexp}
The exponent of the quotient group
\begin{equation}\label{eqn:congexp}
   \frac{S_2(\Gamma; \Z)} { S_2(\Gamma; \Z)[I] + S_2(\Gamma;\Z)[I]^{\perp}}
\end{equation}
is the {\em congruence exponent} $\rAe$ of~$A$ and its
order is the {\em congruence number} $\rA$.
\end{defi}

\begin{rmk}
  Note that $S_2(\Gamma,\Z)\tensor_\Z R = S_2(\Gamma,R)$; see, e.g.,
  the discussion in \cite[\S12]{diamond-im}.  Thus the analogue of
  Definition~\ref{def:congexp} with $\Z$ replaced by an algebraic
  integer ring (or even $\Zbar$) gives a torsion module whose
  annihilator ideal meets~$\Z$ in the ideal generated by the
  congruence exponent.
\end{rmk}

Our definition of~$\rA$ generalizes the definition in
Section~\ref{congintro} when~$A$ is an elliptic curve (see
\cite[p.~276]{abbull}), and the following generalizes 
Theorem~\ref{thm:ribet_au}:
\begin{thm}\label{thm:ribet_gen}
If $f \in S_2(\C)$ is a newform, then
\begin{itemize}
\item[(a)] We have $\nAfe \mid \rAfe$, and
\item[(b)] If $p^2 \nmid N$, then $\ord_p(\rAfe) = \ord_p(\nAfe)$.
\end{itemize}
% $p \nmid \frac{\rAfe}{\nAfe}$.
\end{thm}
%We give the proof of this theorem in the next two sections.
%The rest of the section is devoted to proving Proposition~\ref{ndivsm}
%below, which asserts that if~$f$ is a newform, then $\nAfe \mid
%\rAfe$. 

\begin{rmk}\label{rem:24}
  When $A_f$ is an elliptic curve, Theorem~\ref{thm:ribet_gen} implies
  that the modular degree divides the congruence number (since for an
  elliptic curve the modular degree and modular exponent are the
  same), i.e., $\sqrt{\nAf} \mid \rAf$.  In general, the divisibility
  $\nAf\mid r^2_{A_f}$ need not hold.  For example, there is a newform
  of degree $24$ in $S_2(\Gamma_0(431))$ such that
 $$\nAf = (2^{11}\cdot 6947)^2 \,\,\nmid\,\, r^2_{A_f} = (2^{10}\cdot
  6947)^2.$$  
Note that $431$ is prime and mod~$2$ multiplicity one fails for $J_0(431)$ (see
  \cite{kilford}). 
%The following Magma session illustrates how to verify the above
%assertion about $\nAf$ and $\rAf$.  The commands are parts of Magma
%V2.11 or greater. \vspace{-1ex}
%{\small
%\begin{verbatim}
%       > A := ModularSymbols("431F");
 %      > Factorization(ModularDegree(A));
%       [ <2, 11>, <6947, 1> ]
%       > Factorization(CongruenceModulus(A));
%       [ <2, 10>, <6947, 1> ]
%\end{verbatim}
%}
\end{rmk}


\section{Proof of the Main Theorem}\label{sec:main}
In this section we prove Theorem~\ref{thm:ribet_gen}.
We continue using the notation introduced so far.

\subsection{Proof of Theorem~\ref{thm:ribet_gen} (a)}
\label{sec:firstpart}

We begin with a remark about compatibilities.  In general, the
polarization of~$J$ induced by the theta divisor need not be Hecke
equivariant, because if~$T$ is a Hecke operator on~$J$, then
on~$\Jdual$ it acts as $W_N T W_N$, where $W_N$ is the Atkin-Lehner
involution (see e.g.,~\cite[Rem.~10.2.2]{diamond-im}).  However,
on~$J^{\rm new}$ the action of the Hecke operators commutes with that
of~$W_N$, so if the quotient map $J \ra A$ factors through~$J^{\rm new}$,
then the Hecke action on~$\Adual$ induced by the embedding $\Adual \to
J^{\vee}$ and the action on $\Adual$ induced by $\phi_1:\Adual\to{}J$
are the same.  Hence $\Adual$ is isomorphic to $\po(\Adual)$ 
as a $\T$-module.

Recall that $f$ is a newform, $I_f = {\rm Ann}_\T (f)$, and
$J=J_0(N)$.  Let $B = I_fJ$, so that $\Adual+B=J$, and $J/B\isom A$.
The following lemma is well known, but we prove it here for the
convenience of the reader.

\begin{lem}\label{lem:homzero}
$\Hom_\Q(\Adual,B)=0$. 
\end{lem}
\begin{proof}
%   Suppose there were a nonzero element of $\Hom_\Q(\Adual,B)$.  Since
%   $A$ is simple, for all~$\ell$ the Tate module
%   $V_{\ell}(\Adual)=\Q\tensor\varprojlim_n \Adual[\ell^n]$ would
%   be a factor of $V_{\ell}(B)$.  
% Thus the characteristic polynomial 
% The Eichler-Shimura relation then implies that the characteristic
% polynomial of each 
% One could then extract almost all
%   prime-indexed coefficients of the corresponding eigenforms from the
%   Tate modules, which would violate multiplicity one for systems of
%   Hecke eigenvalues .

  Pick a prime $\ell$.  Then $\Qbar_{\ell} \tensor V_{\ell} (J)^{\ss}$
  as a $\Qbar_{\ell}[G_\Q]$-module is a direct sum of copies of the
  representations $\rho_g$ as $g$ ranges through all normalized
  eigenforms of weight $2$ and level $N$ with coefficients in $\Qbar$;
  by a well-known result of the second author, these representations
  are absolutely irreducible.  Now since~$f$ is a newform and
  $A^{\vee} \to A$ is an isogeny, $\Qbar_{\ell} \tensor
  V_{\ell}(A^{\vee})^{\ss}$ is a direct sum of copies of
  $\rho_{\sigma(f)}$ as $\sigma$ ranges over all embeddings of $K_f$
  into $\Qbar$.  Thus, by the analytic theory of multiplicity one (see
  \cite[Cor.~3, pg.~300]{winnie:newforms}), the Galois modules
  $V_{\ell}(A^{\vee})$ and $V_{\ell}(B) =
  V_{\ell}(J)/V_{\ell}(A^{\vee})$ share no common Jordan-H\"older
  factors even when coefficients are extended to $\Qbar_{\ell}$, so
  $\Hom_\Q(A',B) = 0$.
\end{proof}


Let $\T_1$ be the image of~$\T$ in $\End(\Adual)$, 
and let $\T_2$ be the image of $\T$ in $\End(B)$. 
We have the following commutative diagram with exact rows:
\begin{equation}\label{eqn:diagram}
\xymatrix@=2em{
 0\ar[r] & {\T} \ar[r]\ar[d] & {\T_1\oplus \T_2} \ar[r]\ar[d] & 
                             {\displaystyle \frac{\T_1 \oplus \T_2}{\T}}\ar@{.>}[d]\ar[r] & 0\\
 0\ar[r] & {\End(J)} \ar[r] & {\End(\Adual)\oplus\End(B)} \ar[r] & 
          {\displaystyle \frac{\End(\Adual)\oplus\End(B)}{\End(J)}}\ar[r] & 0.\\
}
\end{equation}
Let 
$$
e=(1,0)\in \T_1 \oplus \T_2,
$$
and let $e_1$ and $e_2$ denote the images of~$e$ in the groups $(\T_1
\oplus \T_2)/\T$ and $(\End(\Adual) \oplus \End(B))/\End(J)$,
respectively.  It follows from Lemma~\ref{lem:homzero} that the two
quotient groups on the right hand side of (\ref{eqn:diagram}) are
finite, so~$e_1$ and~$e_2$ have finite order.  Note that because $e_2$
is the image of $e_1$, the order of $e_2$ is a divisor of the order of
$e_1$.

%this will be used in the proof of Proposition~\ref{ndivsm}
%below.


The {\em denominator} of any $\vphi\in\End(J)\tensor\Q$ is the
smallest positive integer~$n$ such that $n\vphi\in\End(J)$. 
% Explicitly, the denominator of~$\vphi$ is the least common multiples
% of the denominators of the entries of any matrix that represents the
% action of $\vphi$ on the lattice $\H_1(J,\Z)$. 

Let $\piAd, \piB \in \End(J)\tensor\Q$ be projection onto
$\Adual$ and $B$, respectively.  Note that the denominator of
$\piAd$ equals the denominator of $\piB$, since $\piAd
+ \piB = 1_J$, so that $\piB = 1_J - \piAd$. 

\begin{lem}\label{lem:ord_e2}
The element $e_2\in (\End(\Adual) \oplus \End(B))/\End(J)$ 
defined above has order $\nAe$. 
\end{lem}
\begin{proof}
Let $n$ be the order of $e_2$, so~$n$ is the denominator
of $\piAd$, which, as mentioned above, is also the
denominator of $\piB$. We want to show that $n$ is equal
to~$\nAe$, the exponent of $\Adual\cap B$. 

Let $i_{\Adual}$ and $i_B$
be the embeddings of $\Adual$ and $B$ into $J$, respectively. 
Then $$\vphi = (n\piAd,n\piB)\in\Hom(J,\Adual\times B)$$
and $\vphi\circ (i_{\Adual} + i_B) = [n]_{\Adual\times B}.$
We have an exact sequence
$$
0\to \Adual\cap B\xra{x\mapsto (x,-x)}\Adual\times B \xra{i_{\Adual} + i_B}  J \to 0. 
$$
Let $\Delta$ be the image of $\Adual\cap B$.  Then by exactness,
$$
 [n]\Delta =  (\vphi\circ (i_{\Adual} + i_B))(\Delta) = 
\vphi\circ ((i_{\Adual} + i_B)(\Delta)) = \vphi(\{0\}) = \{0\},
$$
so $n$ is a multiple of
the exponent~$\nAe$ of $\Adual\cap B$. 

To show the opposite divisibility, consider the
commutative diagram
$$
\xymatrix@=4em{
0 \ar[r] & {\Adual \cap B} \ar[r]^{x\mapsto (x,-x)}\ar[d]^{[\nAe]}& 
       {\Adual \times B}\ar[d]^{([\nAe],0)}
                  \ar[r]& J \ar[r]\ar@{.>}[d]^{\psi} & 0\\
0 \ar[r] & {\Adual \cap B} \ar[r]^{x\mapsto (x,-x)}& {\Adual \times B}
                  \ar[r]& J \ar[r] & 0,
}
$$
where the middle vertical map is $(a,b)\mapsto (\nAe a,0)$
and the map~$\psi$ exists because $[\nAe](\Adual\cap B)=0$. 
But $\psi = \nAe \piAd$ in $\End(J)\tensor\Q$. 
This shows that $\nAe \piAd \in \End(J)$, i.e.,
that $\nAe$ is a multiple of the
denominator~$n$ of $\piAd$. 

\end{proof}

Let $\Ext^1 = \Ext^1_{\Z}$ denote the first $\Ext$ functor
in the category of $\Z$-modules.

\begin{lem}\label{lem:compare_with_dual}
The group $(\T_1 \oplus \T_2)/\T$ is isomorphic to 
the quotient (\ref{eqn:congexp})
 in Definition~\ref{def:congexp}, so 
 $\rA = \#((\T_1 \oplus \T_2)/\T)$ and $\rAe$ is the
exponent of $(\T_1 \oplus \T_2)/\T$. 
More precisely,  $\Ext^1((\T_1 \oplus \T_2)/\T,\Z)$ is isomorphic as a
$\T$-module to the quotient (\ref{eqn:congexp}). 
\end{lem}
\begin{proof}
Apply the $\Hom(-,\Z)$ functor to the first row of (\ref{eqn:diagram})
to obtain a three-term exact sequence
\begin{equation}\label{eqn:dualseq}
0  \to \Hom(\T_1\oplus \T_2,\Z) \to \Hom(\T,\Z)
\to \Ext^1((\T_1\oplus\T_2)/\T,\Z) \to 0. 
\end{equation}
%The term $\Ext^1(\T_1\oplus \T_2,\Z)$ is $0$ is because
%$\Ext^1(M,\Z)=0$ for any finitely generated free abelian group.  Also,
%$\Hom((\T_1\oplus\T_2)/\T,\Z)=0$ since $(\T_1\oplus\T_2)/\T$ is
%torsion.  
There is a $\T$-equivariant bilinear pairing $\T\times
S_2(\Z)\to\Z$ given by $(t,g)\mapsto a_1(t(g))$, which is perfect by
\cite[Lemma~2.1]{abbull} (see also~\cite[Theorem~2.2]{ribet:modp}). 
Using this pairing, we transform (\ref{eqn:dualseq}) into an exact
sequence
$$
0 \to S_2(\Z)[I_f] \oplus S_2(\Gamma;\Z)[I_f]^{\perp} \to S_2(\Z) \to
\Ext^1((\T_1\oplus\T_2)/\T,\Z) \to 0
$$
of $\T$-modules. 
Here we use that $\Hom(\T_2,\Z)$ is the unique saturated
Hecke-stable complement of $S_2(\Z)[I_f]$ in $S_2(\Z)$, hence
must equal $S_2(\Z)[I_f]^{\perp}$. 
Finally note that if~$G$ is any finite abelian group, then
$\Ext^1(G,\Z)\approx G$ as groups, which gives the desired result. 
\end{proof}

\begin{lem}\label{lem:ord_e1}
The element $e_1 \in (\T_1 \oplus \T_2)/\T$ has order $\rAe$. 
\end{lem}
\begin{proof}
  By Lemma~\ref{lem:compare_with_dual}, the lemma is equivalent to the
  assertion that the order~$r$ of~$e_1$ equals the exponent of
  $M=(\T_1 \oplus \T_2)/\T$.  Since $e_1$ is an element of~$M$, the
  exponent of~$M$ is divisible by~$r$.  
  
  To obtain the reverse divisibility, consider any element $x$ of~$M$.
  Let $(a,b)\in\T_1\oplus \T_2$ be such that its image in~$M$ is~$x$.
  By definition of $e_1$ and~$r$, we have $(r,0)\in\T$, and since
  $1=(1,1)\in\T$, we also have $(0,r)\in\T$.  Thus $(\T{}r,0)$ and
  $(0,\T{}r)$ are both subsets of $\T$ (i.e., in the image of $\T$
  under the map $\T\to\T_1\oplus \T_2$), so $r(a,b)
  =(ra,rb)=(ra,0)+(0,rb)\in \T$.  This implies that the order of~$x$
  divides~$r$. Since this is true for every $x \in M$, we conclude
  that the exponent of~$M$ divides~$r$.
\end{proof}


\begin{prop} \label{ndivsm}
If $f \in S_2(\C)$ is a newform, then
$\nAfe \mid \rAfe$. 
\end{prop}
\begin{proof}
  Since~$e_2$ is the image of~$e_1$ under the right-most
vertical homomorphism in (\ref{eqn:diagram}), the order
  of~$e_2$ divides that of~$e_1$.  Now
apply Lemmas~\ref{lem:ord_e2} and \ref{lem:ord_e1}. 
\end{proof}

This finishes the proof of the first statement in 
Theorem~\ref{thm:ribet_gen}.




\subsection{Proof of Theorem~\ref{thm:ribet_gen} (b)}
\label{sec:secondpart}
Let $\T'$ be the saturation of $\T=\Z[\ldots, T_n,\ldots]$ in
$\End(J_0(N))$, i.e., the set of elements of $\End(J_0(N)) \tensor \Q$
some positive multiple of which lie in~$\T$.
%so 
%$$
% \T' = (\T\tensor\Q) \cap \End(J_0(N)),
%$$
%where the intersection is taken inside $\End(J_0(N))\tensor\Q$.  
The
quotient $\T'/\T$ is a finitely generated abelian group because both
$\T$ and $\End(J_0(N))$ are finitely generated over~$\Z$. Since 
$\T'/\T$ is also a torsion group, it is finite.

In Section~\ref{sec:multone}, we will give some conditions under
which $\T$ and~$\T'$ agree locally at  maximal ideal of~$\T$.
In Section~\ref{sec:degrees}, we will explain how the ratio of
the congruence number to the modular degree is closely related
to the order of~$\T'/\T$, and finally deduce that this ratio is $1$ 
(for quotients associated to newforms) locally at a prime~$p$
such that $p^2 \nmid N$.

\subsubsection{Multiplicity One} \label{sec:multone}

%Suppose for the moment that $M=1$, so $p=pM$.  
Fixt an integer $N$ and a prime $p\mid N$.
Suppose for a moment that $N$ is prime, so $p=N$.
In \cite{mazur:eisenstein},
Mazur proves that $\T=\T'$; he combines this result with
the equality 
$$
  \T\tensor\Q = \End(J_0(p)) \tensor\Q,
$$ 
to deduce that $\T=\End(J_0(p))$.
This result, combined with Ribet's result 
\cite{ribet:endo} or \cite{ribet:endalg}
to the effect that $\T\tensor\Q = (\End_{\Qbar} J_0(N)) \tensor \Q$,
shows that~$\T$ is the full ring of endomorphisms of $J_0(N)$ over $\Qbar$.
When $N$ is no
longer necessarily prime,
the method of \cite{mazur:eisenstein} shows
that $\T$ and $\T'$ agree locally at a maximal ideal $\m$  of $\T$ 
that satisfies a simple condition involving 
differentials form mod $\ell$, where $\ell$ is the residue
characteristic of $\m$.  
%has dimension at most one.
\comment{
$\Supp_{\T}(\T'/\T)$ contains no maximal ideal $\m$ of~$\T$
for which his space $\H^0(X_0(pM)_{\Fell},\Omega)[\m]$ has
dimension $\leq 1$.  (Here $\ell$ is the residue characteristic
of $\m$.)  In other words, multiplicity one for 
$\H^0(X_0(pM)_{\Fell}, \Omega)[\m]$ implies that $\T$ and
$\T'$ agree at~$\m$.
Mazur's argument (see \cite[pg.~95]{mazur:eisenstein}) is quite
general; it relies on a multiplicity $1$ statement for spaces
of differentials in positive characteristic (see 
\cite[Prop.~9.3, pg.~94]{mazur:eisenstein}).  
}

For the sake of completeness, we state and prove a lemma
that can be easily extracted from~\cite{mazur:eisenstein}.  
Let $m$ be the largest square dividing $N$ and
let $R = \Z[\frac{1}{m}]$. Let $X_0(N)_{R}$ denote
the minimal regular model of $X_0(N)$ over~$R$.
%Let $\m$ be a maximal ideal of the Hecke algebra of
%residue characteristic~$\ell$ and suppose $\ell^2 \nmid N$. 
Let $\Omega=\Omega_{X_0(N)/R}$ denote the sheaf of regular
differentials on $X_0(N)_{R}$, as in~\cite[\S2(e)]{mazur:rational}.
If~$\ell$ is a prime such that $\ell^2 \nmid N$, then
$X_0(N)_{\F_\ell}$ denotes the special fiber of $X_0(N)_{R}$ at the
prime~$\ell$.

\comment{
His method shows in 
the general case (where~$M$ is no longer constrained to be~$1$)
that $\Supp_{\T}(\T'/\T)$ contains no maximal ideal $\m$ of~$\T$
for which his space $\H^0(X_0(pM)_{\Fell},\Omega)[\m]$ has
dimension $\leq 1$.  (Here $\ell$ is the residue characteristic
of $\m$.)  In other words, multiplicity one for 
$\H^0(X_0(pM)_{\Fell}, \Omega)[\m]$ implies that $\T$ and
$\T'$ agree at~$\m$.  We record this fact as a lemma
(see also Section~\ref{sec:dataind} for related data).
}

\begin{lem}[Mazur]\label{lem:m1}
Let $\m$ be a maximal ideal of $\T$ of residue characteristic~$\ell$
such that $\ell^2 \nmid N$. 
Suppose that 
$$
 \dim_{\T/\m} \H^0(X_0(N)_{\Fell},\Omega)[\m] \leq 1.
$$
Then $\T$ and $\T'$ agree locally at~$\m$.
%$\m$ is not in the support of $\T'/\T$.
\end{lem}

\begin{proof}
Let $M$ denote the group
$H^1(X_0(N)_R, \OO_{X_0(N)})$,
where $\OO_{X_0(N)}$ is the structure sheaf of $X_0(N)$.
As explained in~\cite[p.~95]{mazur:eisenstein}, 
we have an action of $\EJ$ on~$M$, and 
the action of $\T$ on~$M$ via the inclusion $\T  \subseteq \EJ$
is faithful, so likewise for the action by $\T'$. Hence we have an injection
$\phi: \T' \hookrightarrow {\rm End}_{\T} M$.
% is a free module over~$\T$.
Suppose~$\m$ is a maximal ideal of~$\T$ that satisfies the hypotheses
of the lemma. 
To prove that $\T_\m=\T_\m'$ it suffices to
prove the following claim:\\

\noindent {\em Claim:} The map~$\phi|_{\T}$ is surjective locally at~$\m$.
\begin{proof}
By Nakayama's lemma, to show that $M$ is generated
as a single element over~$\T$ locally at~$\m$,
%to show that $\phi$ is surjective, i.e., 
%to show that $M \tensor \T_\m$ is generated by a single element over
%$\T\tensor \T_\m$, 
it suffices to check that the dimension of the ${\T/\m}\,$-vector space
$M / \m M$ is at most one.
% for each maximal ideal~$\m$ dividing~$\ell$ in~$\T$.  
Since \mbox{$\ell^2 \nmid N$}, 
%$H^1(X_0(N)_{\F_\ell}, \OO) / \m H^1(X_0(N)_{\F_\ell}, \OO)$ 
$M/ \m M$ is dual to  
$H^0(X_0(N)_{\F_\ell}, \Omega) [\m]$ (see, e.g.,~\cite[\S2]{mazur:rational}).
Since we are assuming that 
${\rm dim}_{\T/\m} H^0(X_0(N)_{\F_\ell}, \Omega) [\m] \leq 1$, we have
${\rm dim}_{\T/\m} (M/ \m M) \leq 1$, which proves the claim.
\end{proof}

%We shall use the subscript~$(\m)$ to denote localization at~$\m$.
%Thus $\Mm$ is free of rank one over~$\Tm$. The
%composite $\psi: \Tpm \ra {\rm End}_{\Tm} (\Mm)
%\stackrel{{\phi}^{-1}}{\ra} \Tm$ gives a section of the inclusion
%$\Tm \hookrightarrow \Tpm$. 
%Let $x \in \Tpm$, and let $n$ be an integer such that $nx \in \Tm$.
%Let $y = \psi(x) \in \Tm$. Then $nx = \psi ( \phi (nx) ) = \psi(nx)
%= n \psi(x) = ny$. Since $\Tm$ is torsion-free, this means that
%$x = y \in \Tm$. Thus $\Tm = \Tpm$, as was to be shown.

\comment{
Thus $M \tensor \T_\m$ is free of rank one over~$\T_\m$. The
composite $\EJ \tensor \T_\m
\ra {\rm End}_{\T_\m} (M \tensor \T_\m)
\stackrel{{\phi}^{-1}}{\ra} \T_\m$ gives a section of the inclusion
$\T_\m \hookrightarrow \EJ \tensor \T_\m$. This shows
that $\T_\m$ is saturated in $\EJ \tensor \T_\m$, i.e.,
that $\T$ and $\T'$ agree locally at~$\m$.
}
\end{proof}

If $\m$ is a maximal ideal of the Hecke algebra~$\T$ 
of residue characteristic~$\ell$, we say that
$\m$ satisfies {\em multiplicity one for differentials} if 
$$
\dim (\H^0(X_0(N)_{\F_\ell},\Omega)[\m]) \leq 1.
$$
By Lemma~\ref{lem:m1}, multiplicity one for 
$\H^0(X_0(N)_{\Fell}, \Omega)[\m]$ implies that $\T$ and
$\T'$ agree at~$\m$.

There is quite a bit of literature on the question of multiplicity~$1$
for $\H^0(X_0(N)_{\Fell},\Omega)[\m]$.  
The easiest case is that~$\ell$ is prime to the level $N$: 

\begin{lem}[Mazur]\label{lem:m_ell}
If $\m$ is a maximal ideal of $\T$ of residue characteristic~$\ell$
such that $\ell \nmid N$, then 
$$\dim_{\T/\m} \H^0(X_0(N)_{\Fell},\Omega)[\m] \leq 1.$$ 
%If $\ell \nmid pM$, then~$\ell\nmid \#(\T'/\T)$.
\end{lem}

\begin{proof} 
  Mazur deduces this lemma from injectivity of the $q$-expansion map.
  The reader may find the following alternative approach to part of
  the argument easier to follow than the one on p.~95 of
  \cite{mazur:eisenstein}.  We have an $\Fell$-vector space that
  embeds in $\Fell[[q]]$, for example a space~$V$ of differentials
  that is killed by a maximal ideal $\m$.  This space is a
  $\T/\m$-vector space, and we want to see that its dimension over
  $\T/\m$ is at most~$1$.  Mazur invokes tensor products and
  eigenvectors; alternatively, we note that~$V$ embeds in
  $\Hom_{\Fell}(\T/\m,\Fell)$ via the standard duality that
  sends~$v\in V$ to the linear form whose value on a Hecke
  operator~$T$ is the $q$th coefficient of $v{|T}$.  The group
  $\Hom_{\Fell}(\T/\m,\Fell)$ has the same size as $\T/\m$, which
  completes the argument because $\Hom_{\Fell}(\T/\m,\Fell)$ has
  dimension $1$ as a $\T/\m$-vector space.
\end{proof}
% proves that
%$$\dim_{\T/\m} \H^0(X_0(pM)_{\Fell},\Omega)[\m] \leq 1$$ for 
%all $\m\mid \ell$.    Now apply Lemma~\ref{lem:m1}
%\edit{Is there a problem if $\ell=2$?  How do Lloyd Kilford's examples
%fit into this, where I guess $N=1$ and $\ell=2$ and multiplicity
%one in $J_0(p)$ fails.  Is it still OK in Mazur's
%differentials? -WAS}

In the context of Mazur's paper, where the level~$N$ is prime, we see
from Lemma~\ref{lem:m_ell} that $\T$ and $\T'$ agree away from~$N$.
Locally at $N$, Mazur proved that $\T=\T'$ by an analogue of the
arguments that he used away from $N$; see Chapter II of
\cite{mazur:eisenstein} (and especially Prop.~9.4 and 9.5 of that
chapter) as well as \cite{mazur-ribet}, where these arguments are
taken up in a context where the level is no longer necessarily prime
(and where one works locally at a prime whose square does not divide
the level).  
%At~$N$, we can still use the $q$-expansion principle
%because of the arguments in \cite[Ch.II~\S4]{mazur:eisenstein}.  
Thus in the prime level case, $\T=\T'$, as we asserted above.


Now\label{NPnotation} 
let $p$ be a prime such that $p \parallel N$, and let $M = N/p$.
The question of multiplicity $1$ at $p$ for $\H^0(X_0(pM)_{\Fp},
\Omega)[\m]$ is discussed in \cite{mazur-ribet}, where the authors
establish multiplicity~$1$ for maximal ideals $\m\mid p$ for which the
associated mod~$p$ Galois representation is irreducible and {\em not}
$p$-old.  (A representation of level $pM$ is $p$-old if it arises from
$S_2(\Gamma_0(M))$.)

If~$\m$ is a maximal ideal of~$\T$ of residue characteristic~$\ell$, 
then we say that~$\m$ is ordinary
if $T_\ell \not\in \m$ (note that $T_\ell$ is often denoted $U_\ell$ 
if $\ell \mid N$). For our purposes, the following lemma is convenient:

\begin{lem}[Wiles]\label{lem:wiles}
If $\m$ is an ordinary maximal ideal of $\T$ of characteristic~$p$, then
$$
 \dim_{\T/\m} \H^0(X_0(pM)_{\Fp},\Omega)[\m] \leq 1.
$$
%and $\ord_{\ell}(pM)=1$, then $\m$ is not in the support of $\T'/\T$.
\end{lem}

This is essentially Lemma~2.2 in~\cite[pg.~485]{wiles};
\comment{, which
proves, under a suitable hypothesis, that $\H^0(X_0(pM)_{\F_p},\Omega)[\m]$
is $1$-dimensional if $\m$ is a maximal ideal of~$\T$ that divides~$p$.
The ``suitable hypothesis'' is that $\m$ is ordinary, in the sense that
$T_p \not\in\m$.  (Note that $T_p$ is often denoted $U_p$ in this context.)
It follows from Wiles's lemma that $\T'=\T$ locally at~$\m$ whenever
$\m$ is an ordinary prime whose residue characteristic exactly
divides the level (which is $pM$ here).
}
we make a few comments about how it applies on our situation:
\begin{enumerate}
\item 
Wiles considers $X_1(M,p)$ instead of $X_0(pM)$, which means that he is
using $\Gamma_1(M)$-structure instead of $\Gamma_0(M)$-structure.
This surely has no relevance to the issue at hand.  

\item Wiles assumes (on page 480) that $p$ is an odd prime, but again
this assumption is not relevant to our question.  

\item 
The condition that $\m$ is ordinary does not appear
explicitly in the statement of Lemma~2.2 in~\cite{wiles}; 
instead it is a reigning
assumption in the context of his discussion.  

\item We see by example that Wiles's ``ordinary'' assumption is less
  stringent than the assumption in \cite{mazur-ribet}; note that
  \cite{mazur-ribet} rule out cases where $\m$ is both old and new at
  $p$, whereas Wiles is happy to include such cases.  (On the other
  hand, Wiles's assumption is certainly nonempty, since it rules out
  maximal ideals $\m$ that arise from non-ordinary (old) forms of
  level~$M$.  Here is an example with $p=2$ and $M=11$, so $N=22$:
  There is a unique newform $f=\sum a_n q^n$ of level~$11$, and
  $\T=\Z[T_2] \subset \End(J_0(22))$, where $T_2^2-a_2 T_2 + 2 =0$.
  Since $a_2=-2$, we have $\T\isom \Z[\sqrt{-1}]$.  We can choose the
  square root of $-1$ to be $T_2+1$.  Then $T_2$ is a generator of the
  unique maximal ideal $\m$ of $\T$ with residue characteristic~$2$,
  and this maximal ideal is not ordinary.)
\end{enumerate}
%\end{proof}

We now summarize the conclusions we can make from the lemmas so far:

\begin{prop} \label{prop:TT'}
The modules~$\T$ 
and $\T'$ agree locally at each maximal ideal~$\m$ that is either prime
to~$N$ or that satisfies the following supplemental hypothesis: the
residue characteristic of~$\m$ divides~$N$ only to the first power
and $\m$ is ordinary.  
\end{prop}
\begin{proof}
This follows easily from Lemmas~\ref{lem:m1}, \ref{lem:m_ell},
and~\ref{lem:wiles}.
\end{proof}
\comment{
Wiles's lemma and the standard $q$-expansion argument
(Lemma~\ref{lem:m_ell} and Lemma~\ref{lem:wiles}) imply that~$\T$
and~$\T'$ agree locally at each rational prime that is prime to the
level $pM$, and also at each maximal ideal~$\m$ dividing~$p$ that is
ordinary, in the sense that $T_p \not\in \m$.  A more palatable
description of the situation involves considering the Hecke
algebra~$\T$ and its saturation~$\T'$ at some level $N\geq 1$.  Then
$\T=\T'$ locally at each maximal ideal $\m$ that is either prime
to~$N$ or that satisfies the following supplemental hypothesis: the
residue characteristic of~$\m$ divides~$N$ only to the first power
and~$\m$ is ordinary. 
}

In Mazur's original context, where the level~$N$ is
prime, we have $T_N^2=1$ because there are no forms of
level~$1$. Accordingly, each~$\m$ dividing~$N$ is ordinary, and we
recover Mazur's equality $\T=\T'$ in this special case.

\subsubsection{Degrees and Congruences} \label{sec:degrees}

%Let $e\in \T\tensor\Q$ be an idempotent, and let $A\subset J_0(pM)$
%be the abelian variety image of $e$, i.e., the image of the homomorphism
%$ne\in \T$, where the integer $n\geq 1$ is a multiple of the denominator of $e$.
%Let~$B$ be the image of the complementary idempotent $1-e$.  
%Then $J_0(pM)=A+B$, and $A\cap B$ is a finite group whose exponent
%divides the denominator of $e$.  
%\edit{We can just say that $e$ is as in the previous section.
%Note that $A$ was~$\Adual$ in the previous section. --Amod}

Let $e \in \T\tensor\Q$ be as in Section~\ref{sec:firstpart},
and let $p,N,M$ be as before Lemma~\ref{lem:wiles}.
The image of $e$ in $J_0(pM)$ is the $\T$-stable abelian subvariety
denoted $\Adual$ in Section~\ref{sec:firstpart}, but since we shall
now exclusively work with this subvariety rather than the
corresponding optimal quotient of $J_0(pM)$ (which was denoted $A$
earlier), we will now write $A$ to denote the image of $e$ (without
risk of confusion).  We also write $B$ to denote the unique
$\T$-stable abelian subvariety of $J_0(pM)$ complementary to~$A$.

For $t \in \T$, let $t_A$ be the restriction of~$t$ to $A$, and
let~$t_B$ be the image of~$t$ in $\End(B)$.  Let $\T_A$ be the
subgroup of $\End(A)$ consisting of the various $t_A$, and define
$\T_B$ similarly.  As before, we obtain an injection
$
  j : \T \hra \T_A \times \T_B
$
with finite cokernel.  Because~$j$ is an injection, we
refer to the maps $\pi_A:\T\to \T_A$ and $\pi_B : \T \to \T_B$,
given by $t \mapsto t_A$ and $t\mapsto t_B$, respectively, 
as ``projections''. 

\begin{defi}
The {\em congruence ideal} associated with the projector~$e$ is 
$I=\pi_A(\ker(\pi_B)) \subset \T_A.$
\end{defi}

Viewing $\T_A$ as $\T_A\times \{0\}$, we may view $\T_A$ as a subgroup
of $\T\tensor\Q \isom (\T_A\times \T_B)\tensor\Q$.  Also, we may view
$\T$ as embedded in $\T_A\times \T_B$, via the map~$j$.
\begin{lem}\label{lem:i_int}
We have $I=\T_A\cap \T$.
\end{lem}

A larger ideal of $\T_A$ is 
$
  J = \Ann_{\T_A}(A \cap B);
$
it consists of restrictions to $A$ of Hecke operators that
vanish on $A\cap B$.  

\begin{lem}
We have $I\subset J$.
\end{lem}
\begin{proof}
The image in $\T_A$ of an operator that vanishes on $B$ also
vanishes on $A\cap B$.
\end{proof}

\begin{lem}\label{lem:j_int}
We have 
$J = \T_A \cap \End(J_0(pM)) = \T_A \cap \T'.$
\end{lem}
\begin{proof}
This is elementary; it is an analogue of Lemma~\ref{lem:i_int}.
\end{proof}

\begin{prop}\label{prop:ji_inc}
There is a natural inclusion
$
  J/I \hra \T'/\T
$
of $\T$-modules.
\end{prop}
\begin{proof}
Consider the map $\T\to \T\tensor\Q$ given by $t\mapsto te$.
This homomorphism factors through $\T_A$ and yields an injection
$\iota_A : \T_A \hra \T\tensor\Q$.  Symmetrically, we also
obtain $\iota_B : \T_B \hra \T\tensor\Q$.  The
map
$(t_A, t_B)  \mapsto \iota_A(t_A) + \iota_B(t_B)$
is an injection
 $\T_A\times \T_B \hra \T\tensor\Q$.
The composite of this map with the inclusion $j:\T\hra \T_A\times \T_B$
defined above is the natural map $\T\hra \T\tensor\Q$.  We thus have
a sequence of inclusions
$$
  \T \hra \T_A \times \T_B \hra \T\tensor \Q 
        \subset \End(J_0(pM))\tensor\Q.
$$
By Lemma~\ref{lem:i_int} and Lemma~\ref{lem:j_int},
we have $I=\T_A\cap \T$ and $J=\T_A\cap \T'$.
Thus $I=J\cap \T$, where the intersection is taken
inside $\T'$.  Thus
$$
  J/I = J/(J\cap \T) \isom (J+\T)/\T \hra \T'/\T.
$$
\end{proof}

\begin{cor}\label{cor:ji_inc}
If $\m$ is a maximal ideal not in $\Supp_{\T}(\T'/\T)$,
then $\m$ is not in the support of $J/I$, i.e., 
if $\T$ and $\T'$ agree locally at $\m$, then
$I$ and $J$ also agree locally at $\m$.
\end{cor}

Note that the Hecke algebra $\T$ acts on $J/I$ through 
its quotient $\T_A$, 
since the action of~$\T$ on~$I$ and on~$J$ factors through 
this quotient.

Now we specialize to the case where $A$ is ordinary at $p$,
in the sense that the image of $T_p$ in $\T_A$, which we
denote $T_{p,A}$, is invertible modulo every maximal ideal
of $\T_A$ that divides~$p$.  (This case occurs when~$A$ is
a subvariety of the $p$-new subvariety of $J_0(pM)$, since
the square of $T_{p,A}$ is the identity.)  


If $\m\mid p$
is a maximal ideal of $\T$ that arises by pullback from
a maximal ideal of $\T_A$, then~$\m$ is ordinary in the
sense used above.  When $A$ is ordinary at~$p$, it follows
from Proposition~\ref{prop:TT'} and Corollary~\ref{cor:ji_inc}
that $I=J$ locally at~$p$.  The reason is simple: regarding~$I$
and~$J$ as $\T_A$-modules, we realize that we need to test
that $I=J$ at maximal ideals of $\T_A$ that divide~$p$.
These ideals correspond to maximal ideals $\m\mid p$
of $\T$ that are automatically ordinary, so we have $I=J$
locally at $\m$ because of Proposition~\ref{prop:TT'}.
By Proposition~\ref{prop:TT'}, 
we have $\T=\T'$ locally at primes away from the
level $pM$.  Thus we conclude that $I=J$
locally at all primes $\ell\nmid pM$ and also at~$p$,
a prime that divides the level $pM$ exactly once.

Suppose, finally, that $A$ is the abelian variety associated to a
newform~$f$ of level~$pM$.  
%We then have $\T_A=\Z$.  
The ideal $I\subset \T_A$ measures congruences between~$f$ and the space of forms
in $S_2(\Gamma_0(pM))$ that are orthogonal to the space generated
by~$f$.  Also, $A\cap B$ is the kernel in~$A$ of the map
``multiplication by the modular element~$e$''.  
In this case, the inclusion $I\subset J$ corresponds to the divisibility
$
  \tilde{n}_A \mid \tilde{r}_A,
$
and we have equality at primes at which $I=J$ locally.
We conclude that the congruence exponent and the modular exponent
agree both at~$p$ and at primes not dividing $pM$, which completes our
proof of Theorem~\ref{thm:ribet_gen}(b).

\begin{rmk}
The ring
$$
 R = \End(J_0(pM)) \cap (\T_A \times \T_B)
$$
is often of interest, where the intersection is taken
in $\End(J_0(pM))\tensor \Q$.  We proved above that there
is a natural inclusion $J/I \hra \T'/\T$.  This 
inclusion yields an isomorphism 
$
  J/I \xra{\sim} R/\T.
$
Indeed, if $(t_A, u_B)$ is an endomorphism of $J_0(pM)$,
where $t,u \in \T$, then 
$(t_A, u_B) - u  = (t_A, 0)$ is an element of~$J$.
The ideals~$I$ and~$J$ are equal to the extent that the
rings~$\T$ and $R$ coincide.  Even when $\T'$ is bigger than~$\T$,
its subring $R$ may be not far from~$\T$.
\end{rmk}

\section{Failure of Multiplicity One}\label{sec:mult1}
In this section, we discuss examples of failure of multiplicity one
(in two different but related senses). The notion of multiplicity one,
originally due to Mazur~\cite{mazur:eisenstein}, has played an
important role in several places (e.g., in Wiles's proof of Fermat's
last theorem~\cite{wiles}). This notion is closely related to
Gorensteinness of certain Hecke algebras (e.g., see~\cite{tilouine:hecke}).
Kilford~\cite{kilford} found examples of failure of Gorensteinness
(and multiplicity one) at the prime~$2$ for certain prime levels.
Motivated by the arguments in Section~\ref{sec:main}, in this section
we give examples of failure of multiplicity one for primes (including
odd primes) whose square divides the level.

\subsection{Multiplicity One for
  Differentials}\label{sec:dataind}
In connection with the arguments in Section~\ref{sec:main}, especially
Lemmas~\ref{lem:m1} and \ref{lem:wiles}, it is of interest
to compute the index $[\T':\T]$ for various $N$.
We can compute this index in Magma, e.g., the following
commands compute the index for $N=54$:
``{\tt J := JZero(54); T := HeckeAlgebra(J); Index(Saturation(T), T);}''
We obtain Table~\ref{table:index}, where the first column
contains $N$ and the second column contains $[\T':\T]$:
\begin{table}\caption{The Index $[\T':\T]$\label{table:index}}
\begin{center}
\begin{tabular}{|l|c|}\hline
11 & 1 \\\hline
12 & 1 \\\hline
13 & 1 \\\hline
14 & 1 \\\hline
15 & 1 \\\hline
16 & 1 \\\hline
17 & 1 \\\hline
18 & 1 \\\hline
19 & 1 \\\hline
20 & 1 \\\hline
21 & 1 \\\hline
22 & 1 \\\hline
23 & 1 \\\hline
24 & 1 \\\hline
25 & 1 \\\hline
26 & 1 \\\hline
27 & 1 \\\hline
28 & 1 \\\hline
29 & 1 \\\hline
30 & 1 \\\hline
31 & 1 \\\hline
32 & 1 \\\hline
33 & 1 \\\hline
34 & 1 \\\hline
35 & 1 \\\hline
36 & 1 \\\hline
37 & 1 \\\hline
38 & 1 \\\hline
39 & 1 \\\hline
40 & 1 \\\hline
41 & 1 \\\hline
42 & 1 \\\hline
43 & 1 \\\hline
44 & 2 \\\hline
45 & 1 \\\hline
46 & 2 \\\hline
47 & 1 \\\hline
48 & 1 \\\hline
49 & 1 \\\hline
50 & 1 \\\hline
\end{tabular}
\,\,\,\,\,\,
\begin{tabular}{|l|c|}\hline
51 & 1 \\\hline
52 & 1 \\\hline
53 & 1 \\\hline
54 & 3 \\\hline
55 & 1 \\\hline
56 & 2 \\\hline
57 & 1 \\\hline
58 & 1 \\\hline
59 & 1 \\\hline
60 & 2 \\\hline
61 & 1 \\\hline
62 & 2 \\\hline
63 & 1 \\\hline
64 & 2 \\\hline
65 & 1 \\\hline
66 & 1 \\\hline
67 & 1 \\\hline
68 & 2 \\\hline
69 & 1 \\\hline
70 & 1 \\\hline
71 & 1 \\\hline
72 & 2 \\\hline
73 & 1 \\\hline
74 & 1 \\\hline
75 & 1 \\\hline
76 & 2 \\\hline
77 & 1 \\\hline
78 & 2 \\\hline
79 & 1 \\\hline
80 & 4 \\\hline
81 & 1 \\\hline
82 & 1 \\\hline
83 & 1 \\\hline
84 & 2 \\\hline
85 & 1 \\\hline
86 & 1 \\\hline
87 & 1 \\\hline
88 & 8 \\\hline
89 & 1 \\\hline
90 & 1 \\\hline
\end{tabular}
\,\,\,\,\,\,
\begin{tabular}{|l|c|}\hline
91 & 1 \\\hline
92 & 16 \\\hline
93 & 1 \\\hline
94 & 4 \\\hline
95 & 1 \\\hline
96 & 8 \\\hline
97 & 1 \\\hline
98 & 1 \\\hline
99 & 9 \\\hline
100 & 1 \\\hline
101 & 1 \\\hline
102 & 1 \\\hline
103 & 1 \\\hline
104 & 4 \\\hline
105 & 1 \\\hline
106 & 1 \\\hline
107 & 1 \\\hline
108 & 54 \\\hline
109 & 1 \\\hline
110 & 2 \\\hline
111 & 1 \\\hline
112 & 8 \\\hline
113 & 1 \\\hline
114 & 1 \\\hline
115 & 1 \\\hline
116 & 4 \\\hline
117 & 1 \\\hline
118 & 2 \\\hline
119 & 1 \\\hline
120 & 32 \\\hline
121 & 1 \\\hline
122 & 1 \\\hline
123 & 1 \\\hline
124 & 16 \\\hline
125 & 25 \\\hline
126 & 18 \\\hline
127 & 1 \\\hline
128 & 64 \\\hline
129 & 1 \\\hline
130 & 1 \\\hline
\end{tabular}
\,\,\,\,\,\,
\begin{tabular}{|l|c|}\hline
131 & 1 \\\hline
132 & 8 \\\hline
133 & 1 \\\hline
134 & 1 \\\hline
135 & 27 \\\hline
136 & 16 \\\hline
137 & 1 \\\hline
138 & 4 \\\hline
139 & 1 \\\hline
140 & 8 \\\hline
141 & 1 \\\hline
142 & 8 \\\hline
143 & 1 \\\hline
144 & 32 \\\hline
145 & 1 \\\hline
146 & 1 \\\hline
147 & 7 \\\hline
148 & 4 \\\hline
149 & 1 \\\hline
150 & 5 \\\hline
151 & 1 \\\hline
152 & 32 \\\hline
153 & 9 \\\hline
154 & 1 \\\hline
155 & 1 \\\hline
156 & 32 \\\hline
157 & 1 \\\hline
158 & 4 \\\hline
159 & 1 \\\hline
160 & 256 \\\hline
161 & 1 \\\hline
162 & 81 \\\hline
163 & 1 \\\hline
164 & 8 \\\hline
165 & 1 \\\hline
166 & 2 \\\hline
167 & 1 \\\hline
168 & 128 \\\hline
169 & 13 \\\hline
170 & 1 \\\hline
\end{tabular}
\,\,\,\,\,\,
\begin{tabular}{|l|c|}\hline
171 & 9 \\\hline
172 & 8 \\\hline
173 & 1 \\\hline
174 & 4 \\\hline
175 & 5 \\\hline
176 & 512 \\\hline
177 & 1 \\\hline
178 & 1 \\\hline
179 & 1 \\\hline
180 & 72 \\\hline
181 & 1 \\\hline
182 & 1 \\\hline
183 & 1 \\\hline
184 & 1024 \\\hline
185 & 1 \\\hline
186 & 4 \\\hline
187 & 1 \\\hline
188 & 256 \\\hline
189 & 243 \\\hline
190 & 8 \\\hline
191 & 1 \\\hline
192 & 4096 \\\hline
193 & 1 \\\hline
194 & 1 \\\hline
195 & 1 \\\hline
196 & 14 \\\hline
197 & 1 \\\hline
198 & 81 \\\hline
199 & 1 \\\hline
200 & 80 \\\hline
201 & 1 \\\hline
202 & 1 \\\hline
203 & 1 \\\hline
204 & 32 \\\hline
205 & 1 \\\hline
206 & 4 \\\hline
207 & 81 \\\hline
208 & 256 \\\hline
209 & 1 \\\hline
210 & 2 \\\hline
\end{tabular}
\end{center}
\end{table}

Let $\m$ be a maximal ideal of the Hecke algebra
$\T\subset\End(J_0(N))$ of residue characteristic~$p$. Recall
that we say that
$\m$ satisfies {\em multiplicity one for differentials} if $\dim
(\H^0(X_0(N)_{\Fp},\Omega)[\m]) \leq 1$.  

In each case in which $[\T':\T]\neq 1$, Lemma~\ref{lem:m1} implies
that there is some maximal ideal $\m$ of $\T$ such that
$\dim(\H^0(X_0(N)_{\Fp},\Omega)[\m])>1$, which is an example
of failure of multiplicity one for differentials.


In Table~\ref{table:index}, whenever $p\mid [\T':\T]$, then $p^2\mid
2N$.  This is a consequence of Proposition~\ref{prop:TT'}, which
moreover asserts that when $2$ exactly divides $N$ and $2\mid[\T':\T]$
then there is a non-ordinary (old) maximal ideal of characteristic $2$
in the support of $\T'/\T$.

%The first case when $2\mid\mid N$ and $2\mid
%[\T':\T]$ is $N=46$, where we find (via a Magma calculation) that
%$G=\T'/\T \isom \Z/2\Z$, and the Hecke operator~$T_2$ acts as~$0$ on
%$G$, so the annihilator of $G$ in $\T$ is not ordinary, which does not
%ontradict Proposition~\ref{prop:TT'}.

Moreover, notice that
Theorem~\ref{thm:ribet_gen}(b) (whose proof is in 
Section~\ref{sec:secondpart})
follows formally from two
key facts: that $A_f$ is new and that multiplicity one for differentials
holds for ordinary maximal ideals with residue characteristic
$p\mid\mid N$ and for all maximal ideals with residue
characteristic $p\nmid N$.  The conclusion of
Theorem~\ref{thm:ribet_gen}(b) does not hold for the counterexamples
in Section~\ref{sec:elliptic} (e.g., for~54B1), which are
all new elliptic curves, so multiplicity one for
differentials does not hold for certain 
maximal ideals that arise from the new quotient of the Hecke algebra.
Note that in all examples we have $p\mid(r/m)$ with $p^2\mid N$,
which raises the question: are there non-ordinary counterexamples
with $p\mid\mid N$?

%\edit{I think I have been to vague. Basically, I wanted to say
%that one reason why the index $[\T':\T]$ could be nontrivial
%(and hence lead to failure of multiplicity one for differentials)
%is due to old-ness. But even in the new part, one may have failure
%of mult one for diffs, due to the level not being square-free. --Amod}

% was@modular:~/comps/ind_table$ magma
% Magma V2.11-10    Sun Aug 28 2005 18:34:49 on modular  [Seed = 1293693469]
% Type ? for help.  Type <Ctrl>-D to quit.
% > J := JZero(46);
% > J;
% Modular abelian variety JZero(46) of dimension 5 and level 2*23 over Q
% > T := HeckeAlgebra(J);
% > S := Saturation(T);
% > m := S/T;
% > m;
% Abelian Group isomorphic to Z/2
% Defined on 1 generator
% Relations:
%     2*m.1 = 0
% >
% > S;
% Sat(HeckeAlg(JZero(46))): Group of homomorphisms from JZero(46) to JZero(46)
% > Basis(S);
% [
%     Homomorphism from JZero(46) to JZero(46) (not printing 10x10 matrix),
%     Homomorphism from JZero(46) to JZero(46) (not printing 10x10 matrix),
%     Homomorphism from JZero(46) to JZero(46) (not printing 10x10 matrix),
%     Homomorphism from JZero(46) to JZero(46) (not printing 10x10 matrix),
%     Homomorphism from JZero(46) to JZero(46) (not printing 10x10 matrix)
% ]
% > S.1;
% Homomorphism from JZero(46) to JZero(46) (not printing 10x10 matrix)
% > S.1 in T;
% true
% > S.2 in T;
% false
% > t2 := HeckeOperator(J,2);
% > t2 in T;
% true
% > t2*S.2;
% Homomorphism from JZero(46) to JZero(46) (not printing 10x10 matrix)
% > t2*S.2 in T;
% true

\subsection{Multiplicity One for Jacobians}

We say that a maximal ideal~$\m$ of~$\T$ satisfies {\it multiplicity one}
if $J_0(N)[\m]$ is of dimension two over~$\T/\m$. We sometimes use the
phrase ``multiplicity one for~$J_0(N)$'' in order to distinguish this notion
from the notion of multiplicity one for differentials. 
%\edit{I added these lines for clarity. --Amod}

\begin{prop}\label{prop:mult1J}
  Suppose $E$ is an optimal elliptic curve over $\Q$ of conductor~$N$
  and~$p$ is a prime such that $p \mid r_E$ but $p\nmid m_E$.  Let
  $\m$ be the annihilator in $\T$ of $E[p]$.  Then multiplicity one
  fails for $\m$, i.e., $\dim_{\T/\m} J_0(N)[\m] > 2$.
%\edit{Earlier the conclusion said $\dim_{\T/\m} J_0(N)[\m] > 1$.
%--Amod}
\end{prop}
\begin{proof}
  Using the principal polarization $E \isom E^{\vee}$ we view~$E$ as
  an abelian subvariety of $J=J_0(N)$ and consider the complementary
  $\T$-stable abelian subvariety $A$ of $E$ (thus $A$ is the kernel of
  the modular parametrization map $J\to E$).  In this setup, $J = E +
  A$, and the intersection of $E$ and $A$ is $E[m_E]$.  Here we use
  that the composite map
  $
    E \simeq E^{\vee} \to J^{\vee} \to J \to E
  $
  is a polarization, and hence is multiplication by a positive integer
  $m_E$.  Because $p\nmid m_E$, we have $E[p]\cap A = 0$.  On the
  other hand, let $\m$ be the annihilator of $E[p]$ inside $\T$.  Then
  $J[\m]$ contains $E[p]$ and also $A[\m]$, and because $p$ is a
  congruence prime, the submodule $A[\m]\subset J[\m]$ is nonzero.
  Thus the sum $E[p] + A[\m]$ is a direct sum and is larger than
  $E[p]$, which is of dimension $2$ over $\T/\m = \Z/p\Z$.  Hence the
  dimension of $J[\m]$ over $\T/\m$ is bigger than $2$, as claimed.
%\edit{I changed the last two lines. --Amod}
\end{proof}

Proposition~\ref{prop:mult1J} implies that any example in which
simultaneously $p\nmid m_E$ and $\ord_p(r_E)\neq \ord_p(m_E)$ produces
an example in which multiplicity one for $J_0(N)$ fails.  For example,
for the curve 54B1 and $p=3$, we have $\ord_3(r_E)=1$ but
$\ord_3(m_E)=0$, so multiplicity one at $3$ fails for $J_0(54)$.
%Also, for 242B1, we have $r_E = 11\cdot 2^4$ and $m_E = 2^4$, so
%multiplicity one for $J_0(242)$ fails at $11$.  For $N=242$ we also
%have $[\T':\T]=121$, so multiplicity one at $11$ also fails for
%differentials 
%\edit{Doesn't this work for 54B1 as well? Why give another
%example? --Amod}
%(see Section~\ref{sec:dataind} above).

%\edit{I commented out the old section, which had mistakes. --Amod}
\comment{
\subsection{Multiplicity one (old section)}
\edit{william: I think this section should be deleted in light
of the above two new sections.}
%\edit{[This is still very rough, and will need to be cleaned up. --Amod]}

Let $\m$ be a maximal ideal of the Hecke algebra
$\T\subset\End(J_0(N))$ of residue characteristic~$p$. We say that
$\m$ satisfies {\em strong multiplicity one} if $\dim
(\H^0(X_0(N)_{\Fp},\Omega)[\m]) = 1$.  
\edit{or should be it be equal to~$1$? Also ``strong'' is just
  something I came up with; may want to change the name.  In fact, I
  may have got strong and weak mixed up, since Tilouine calls 
  above ``weak multiplicity one''. --Amod} 

The proof of Theorem~\ref{thm:ribet_gen}(b) follows formally from two
key facts: that $A_f$ is new and that strong multiplicity one holds
for ordinary maximal ideals if~$p^2 \nmid N$.  The conclusion of
Theorem~\ref{thm:ribet_gen}(b) does not hold for the counterexamples
in Section~\ref{sec:elliptic} (at levels $54$,~$64$, etc.), which are
all new elliptic curves, which shows that strong multiplicity one does
not hold for certain ordinary\edit{William: I totally don't get this.
When $p^2\mid N$ the ideals are {\em NOT} ordinary.  We only got
ordinary in the proof above because $p\mid\mid N$.  For example, for
the curve 54b, we have $a_3=0$, and the ideal is not ordinary.}
maximal ideals for the corresponding
levels (in all of them, $p^2 \mid N$). We record this observation:

\begin{prop}
  There are ordinary maximal ideals for which strong
  multiplicity one fails.\edit{William: But our examples 
are not ordinary!}
\end{prop}

There is another notion of multiplicity one: 
suppose~$\rho_{\m}$, the representation attached to~$\m$,
is absolutely irreducible.
Then one
says that $\m$ satisfies {\em weak multplicity one} if
$J_0(N)[\m]$ is isomorphic to a single copy of~$\rho_{\m}$.
By standard arguments, strong multiplicity one implies
weak multiplicity one when $\overline{\rho}_{\m}$ is absolutely irreducible.
\edit{William: I think this is wrong.
In connection with Remark~\ref{rem:24}, we find (again via a Magma
computation) that $\T'=\T$ when $N=431$.  This was also already known via
work of Mazur \cite{mazur:eisenstein}, and illustrates that
multiplicity one for differentials need not imply multiplicity one for
$J_0(N)$.  So I'm confused.}
%\edit{Will: I just realized that for this implication,
%one also needs the representation~$\rho_{\m}$
%to be absolutely irreducible -- it might be worth checking
%if this holds in our couterexamples}
which in turn implies that the Hecke algebra~$\T$ is Gorenstein.\edit{or
maybe weak multiplicity one is equivalent to Gorensteinness, and both
follow from strong multiplicity one; a good reference might be
Tilouine's article in the FLT conference.}  Wiles proves that strong
multiplicity one holds for ordinary maximal ideals provided $p^2 \nmid
N$, which he used to show the Gorensteinness of certain Hecke
algebras. This Gorensteinness property was a key step in the proof of
Fermat's last theorem. Our finding above shows that the hypothesis
$p^2 \nmid N$ is essential, and thus gives a limit to how far the
standard argument for proving Gorensteinness works.

Note that while in our examples, we know that strong multiplicity one
fails, we do not know if weak multiplicity one or Gorensteinness
fails. It would be interesting to do calculations for the latter, like
were done by Kilford. 
}

%\bibliographystyle{amsalpha}         
%\bibliography{biblio}
%\end{document}

\providecommand{\bysame}{\leavevmode\hbox to3em{\hrulefill}\thinspace}
\providecommand{\MR}{\relax\ifhmode\unskip\space\fi MR }
% \MRhref is called by the amsart/book/proc definition of \MR.
\providecommand{\MRhref}[2]{%
  \href{http://www.ams.org/mathscinet-getitem?mr=#1}{#2}
}
\providecommand{\href}[2]{#2}
\begin{thebibliography}{BCDT01}

\bibitem[ARS06]{ars}
A.~Agashe, K.~Ribet and W.~Stein, 
\emph{The {M}anin {C}onstant}, QJPAM, Coates Volume (2006), to appear. 

\bibitem[AU96]{abbull}
A.~Abbes and E.~Ullmo, \emph{\`{A} propos de la conjecture de {M}anin pour les
  courbes elliptiques modulaires}, Compositio Math. \textbf{103} (1996), no.~3,
  269--286. 

\bibitem[BCP97]{magma}
W.~Bosma, J.~Cannon, and C.~Playoust, \emph{The {M}agma algebra system. {I}.
  {T}he user language}, J. Symbolic Comput. \textbf{24} (1997), no.~3--4,
  235--265, Computational algebra and number theory (London, 1993). 
%\MR{1 484 478}

\bibitem[BCDT01]{breuil-conrad-diamond-taylor}
C.~Breuil, B.~Conrad, F.~Diamond, and R.~Taylor, \emph{On the modularity of
  elliptic curves over {$\bold Q$}: wild 3-adic exercises}, J. Amer. Math. Soc.
  \textbf{14} (2001), no.~4, 843--939 (electronic). %\MR{2002d:11058}

\bibitem[CK04]{cojo-kani}
Alina~Carmen Cojocaru and Ernst Kani, \emph{The modular degree and the
  congruence number of a weight 2 cusp form}, Acta Arith. \textbf{114} (2004),
  no.~2, 159--167. 

\bibitem[Cre97]{cremona:alg}
J.\thinspace{}E. Cremona, \emph{Algorithms for modular elliptic curves}, second
  ed., Cambridge University Press, Cambridge, 1997, \mbox{}\\{\tt
  http://www.maths.nott.ac.uk/personal/jec/book/}.

\bibitem[CSS97]{boston}
G.~Cornell, J.\thinspace{}H. Silverman, and G.~Stevens (eds.), \emph{Modular
  forms and \protect{F}ermat's last theorem (boston,ma, 1995)}, New York,
  Springer-Verlag, 1997, Papers from the Instructional Conference on Number
  Theory and Arithmetic Geometry held at Boston University, Boston, MA, August
  9--18, 1995.


\bibitem[DI95]{diamond-im}
F.~Diamond and J.~Im, \emph{Modular forms and modular curves}, Seminar on
  {F}ermat's {L}ast {T}heorem, Providence, RI, 1995, pp.~39--133.

\bibitem[Fre97]{frey:ternary}
G.~Frey, \emph{On ternary equations of {F}ermat type and relations with
  elliptic curves}, Modular forms and Fermat's last theorem (Boston, MA, 1995),
  Springer, New York, 1997, pp.~527--548. 

\bibitem[FM99]{frey-muller}
G.~Frey and M.~M{\"u}ller, \emph{Arithmetic of modular curves and
  applications}, Algorithmic algebra and number theory (Heidelberg, 1997),
  Springer, Berlin, 1999, pp.~11--48.

\bibitem[Har77]{hartshorne:ag}
R.~Hartshorne, \emph{Algebraic geometry}, Springer-Verlag, New York, 1977,
  Graduate Texts in Mathematics, No. 52. 

\bibitem[Kil02]{kilford}
L.\thinspace{}J.\thinspace{}P. Kilford, 
  \emph{Some non-{G}orenstein {H}ecke algebras attached to
  spaces of modular forms}, J. Number Theory \textbf{97} (2002), no.~1,
  157--164. 

\bibitem[Lan83]{lang:av}
S.~Lang, \emph{Abelian varieties}, Springer-Verlag, New York, 1983, Reprint
  of the 1959 original. 


\bibitem[Li75]{winnie:newforms}
W-C. Li, \emph{Newforms and functional equations}, Math. Ann. \textbf{212}
  (1975), 285--315.

\bibitem[Maz77]{mazur:eisenstein}
B.~Mazur, \emph{Modular curves and the \protect{Eisenstein} ideal}, Inst.
  Hautes \'Etudes Sci. Publ. Math. (1977), no.~47, 33--186 (1978).

\bibitem[Maz78]{mazur:rational}
B.~Mazur, \emph{Rational isogenies of prime degree (with an appendix by {D}.
  {G}oldfeld)}, Invent. Math. \textbf{44} (1978), no.~2, 129--162.

\bibitem[MR91]{mazur-ribet}
B.~Mazur and K.~A. Ribet, \emph{Two-dimensional representations in the
  arithmetic of modular curves}, Ast\'erisque (1991), no.~196--197, 6, 215--255
  (1992), Courbes modulaires et courbes de Shimura (Orsay, 1987/1988).
  %\MR{MR1141460 (93d:11056)}

\bibitem[Mum70]{mumford:av}
D.~Mumford, \emph{Abelian varieties}, Tata Institute of Fundamental Research
  Studies in Mathematics, No. 5, Published for the Tata Institute of
  Fundamental Research, Bombay, 1970. 

\bibitem[Mur99]{murty:congruence}
M.\thinspace{}R. Murty, \emph{Bounds for congruence primes}, Automorphic forms,
  automorphic representations, and arithmetic (Fort Worth, TX, 1996), Amer.
  Math. Soc., Providence, RI, 1999, pp.~177--192. %\MR{2000g:11038}

\bibitem[Rib75]{ribet:endo}
K.\thinspace{}A. Ribet, \emph{Endomorphisms of semi-stable abelian varieties
  over number fields}, Ann. Math. (2) \textbf{101} (1975), 555--562. 
%\MR{51 \#8120}

\bibitem[Rib81]{ribet:endalg}
K.~A. Ribet, \emph{Endomorphism algebras of abelian varieties attached to
  newforms of weight {$2$}}, Seminar on Number Theory, Paris 1979--80, Progr.
  Math., vol.~12, Birkh\"auser Boston, Mass., 1981, pp.~263--276.
 % \MR{82m:10044}

\bibitem[Rib83]{ribet:modp}
K.\thinspace{}A. Ribet, \emph{Mod $p$\ {H}ecke operators and congruences
  between modular forms}, Invent. Math. \textbf{71} (1983), no.~1, 193--205. 


\bibitem[Ste89]{stevens:param}
G.~Stevens, \emph{Stickelberger elements and modular parametrizations of
  elliptic curves}, Invent. Math. \textbf{98} (1989), no.~1, 75--106.


\bibitem[Stu87]{sturm:cong}
J.~Sturm, \emph{On the congruence of modular forms}, Number theory (New York,
  1984--1985), Springer, Berlin, 1987, pp.~275--280.

\bibitem[Til97]{tilouine:hecke}
J.~Tilouine, \emph{Hecke algebras and the {G}orenstein property}, Modular
  forms and Fermat's last theorem (Boston, MA, 1995), Springer, New York, 1997,
  pp.~327--342.


\bibitem[Vat05]{MR2135139}
V.~Vatsal, \emph{Multiplicative subgroups of {$J\sb 0(N)$} and applications to
  elliptic curves}, J. Inst. Math. Jussieu \textbf{4} (2005), no.~2, 281--316.

\bibitem[Wil95]{wiles}
A.\thinspace{}J. Wiles, \emph{Modular elliptic curves and \protect{F}ermat's
  last theorem}, Ann. of Math. (2) \textbf{141} (1995), no.~3, 443--551. 

\bibitem[Zag85]{zagier}
D.~Zagier, \emph{Modular parametrizations of elliptic curves}, Canad. Math. 
 Bull. \textbf{28} (1985), no.~3, 372--384. 

\end{thebibliography}


\end{document}










